% !TeX root=../main.tex
% در این فایل، عنوان پایان‌نامه، مشخصات خود، متن تقدیمی‌، ستایش، سپاس‌گزاری و چکیده پایان‌نامه را به فارسی، وارد کنید.
% توجه داشته باشید که جدول حاوی مشخصات پروژه/پایان‌نامه/رساله و همچنین، مشخصات داخل آن، به طور خودکار، درج می‌شود.
%%%%%%%%%%%%%%%%%%%%%%%%%%%%%%%%%%%%
% دانشگاه خود را وارد کنید
\university{دانشگاه اصفهان}
% پردیس دانشگاهی خود را اگر نیاز است وارد کنید (مثال: فنی، علوم پایه، علوم انسانی و ...)
\college{پردیس دانشکده‌های فنی}
% دانشکده، آموزشکده و یا پژوهشکده  خود را وارد کنید
\faculty{دانشکده مهندسی کامپیوتر}
% گروه آموزشی خود را وارد کنید (در صورت نیاز)
\department{هوش مصنوعی و رباتیک}
% رشته تحصیلی خود را وارد کنید
\subject{هوش مصنوعی و رباتیک}
% گرایش خود را وارد کنید
\field{هوش مصنوعی}
% عنوان پایان‌نامه را وارد کنید
\title{
‫ﺗﻮﺳﻌﻪ ‬‫ﯾﮏ ‬‫ﻣﺪل‬ ‫زﺑﺎﻧﯽ‬ ‫ﭘ‬ﺰﺷﮑﯽ‬ ‫ﻣﺒﺘﻨﯽ‬ ‫ﺑﺮ‬ ‫اﺳﺘﺪﻻل در زبان فارسی
}
% نام استاد(ان) راهنما را وارد کنید
\firstsupervisor{دکتر حمیدرضا برادران کاشانی}
\firstsupervisorrank{استاد}
% نام داوران داخلی و خارجی خود را وارد نمایید.
\internaljudge{دکتر الهام کردی قصردشتی}
\internaljudgerank{استاد}
\externaljudge{دکتر محمدعلی نعمت بخش}
\externaljudgerank{استاد}
\externaljudgeuniversity{دانشگاه داور خارجی}
% نام نماینده کمیته تحصیلات تکمیلی در دانشکده \ گروه
\graduatedeputy{دکتر فریا نصیری مفخم}
\graduatedeputyrank{استاد}
% نام دانشجو را وارد کنید
\name{مهرداد}
% نام خانوادگی دانشجو را وارد کنید
\surname{قصابی}
% شماره دانشجویی دانشجو را وارد کنید
\studentID{4023614029}
% تاریخ پایان‌نامه را وارد کنید
\thesisdate{بهمن ۱۴۰۴}
% به صورت پیش‌فرض برای پایان‌نامه‌های کارشناسی تا دکترا به ترتیب از عبارات «پروژه»، «پایان‌نامه» و «رساله» استفاده می‌شود؛ اگر  نمی‌پسندید هر عنوانی را که مایلید در دستور زیر قرار داده و آنرا از حالت توضیح خارج کنید.
%\projectLabel{پایان‌نامه}

% به صورت پیش‌فرض برای عناوین مقاطع تحصیلی کارشناسی تا دکترا به ترتیب از عبارت «کارشناسی»، «کارشناسی ارشد» و «دکتری» استفاده می‌شود؛ اگر نمی‌پسندید هر عنوانی را که مایلید در دستور زیر قرار داده و آنرا از حالت توضیح خارج کنید.
%\degree{}
%%%%%%%%%%%%%%%%%%%%%%%%%%%%%%%%%%%%%%%%%%%%%%%%%%%%
%% پایان‌نامه خود را تقدیم کنید! %%
\dedication
{
	{\Large تقدیم به:}\\
	\begin{flushleft}{
			\huge
			هم کلاسی عزیزم محمد جعفرپور\\
		}
	\end{flushleft}
}
%% متن قدردانی %%
%% متن قدردانی %%
%% ترجیحا با توجه به ذوق و سلیقه خود متن قدردانی را تغییر دهید.
\acknowledgement{
 دستان پدر و مادر نازنینم را به پاس مهر بیکرانشان میبوسم و از استاد راهنما خود جناب آقای دکتر حمیدرضا برادران بابت راهنمایی هایشان در طول انجام این پایان نامه سپاس گزاری میکنم.
}
%%%%%%%%%%%%%%%%%%%%%%%%%%%%%%%%%%%%
%چکیده پایان‌نامه را وارد کنید
\fa-abstract{
استفاده از هوش مصنوعی در پاسخگویی به سوالات پزشکی به عنوان یکی از حوزه‌های نوظهور و مهم در فناوری و بهداشت، در سال‌های اخیر توجه زیادی را به خود جلب کرده است. در حالی که تحقیقات گسترده‌ای در این زمینه به زبان انگلیسی منتشر شده، در زمان آغاز این پایان‌نامه، حوزه زبان فارسی خالی از پژوهش بود. در فاز نخست این تحقیق، یک پیکره پزشکی شامل نود میلیون توکن و مجموعه داده‌ای متشکل از بیست هزار پرسش و پاسخ پزشکی که توسط پزشکان پاسخ داده شده، از مجلات و تالارهای پرسش و پاسخ آنلاین جمع‌آوری شد. با استفاده از بخشی از این داده‌ها، مدل پایه «آیا اکسپنس» آموزش داده شد و مدل جدیدی به نام «گائوکرنا وی» معرفی گردید که موفق به کسب
\lr{2.67}
درصد عملکرد بهتر نسبت به مدل پایه خود بر روی مجموعه داده
\lr{MMLU}
پزشکی ترجمه شده به زبان فارسی شد. در فاز دوم، با توجه به اینکه پزشکی به شدت وابسته به استدلال و تحلیل‌های منطقی است، یک مدل جدید مبتنی بر زنجیره‌ای از افکار و استدلال‌های منطقی طراحی شد تا دقت و کارایی آن را به طور قابل توجهی افزایش دهد. برای این منظور، دو چهارچوب جدید بر اساس یادگیری تقویتی با بازخورد هوش مصنوعی معرفی شد تا قابلیت دلیل‌آوری مدل پایه یعنی آیا اکسپنس را از طریق بهینه‌سازی مستقیم ترجیحات بهبود بخشد. نتیجه این بهبود، مدل «گائوکرنا آر» بود که با وجود آموزش بر روی تعداد توکن‌های بسیار کمتر نسبت به گائوکرنا وی، موفق به کسب
\lr{6.34}
درصد عملکرد بهتر نسبت به مدل پایه خود بر روی همان مجموعه داده
\lr{MMLU}
پزشکی ترجمه شده به زبان فارسی شد که این میزان بهبود نشان‌دهنده اهمیت استدلال در حوزه پزشکی است.
}
% کلمات کلیدی پایان‌نامه را وارد کنید
\keywords{هوش مصنوعی در پزشکی، مدل های زبانی فارسی، مدل های زبانی پزشکی،  پردازش زبان های طبیعی، توانایی استدلال هوش مصنوعی}
% انتهای وارد کردن فیلد‌ها
%%%%%%%%%%%%%%%%%%%%%%%%%%%%%%%%%%%%%%%%%%%%%%%%%%%%%%
